% 5-5-conclusion.tex
%  Budget: ~2pp.  Final synthesis. NO new numbers, NO new claims.
%  MANDATE: never "beats the state of the art"; the negative results stay
%  in the closing paragraphs, not buried.
%  Rules: decimals in \lr{}; integers Persian; every table cell in \lr{}.

\section{نتیجه‌گیری نهایی}
\label{sec:5-5}

این پایان‌نامه با پرسشی متعارف آغاز شد \lr{—} اینکه قیمت روزِ بعدِ برق آلمان را
چگونه می‌توان بهتر پیش‌بینی کرد \lr{—} و با پاسخی پایان می‌یابد که بخشی از آن
مثبت است و بخشی نه. هر دو بخش گزارش شده‌اند، و همین انتخاب، بیش از هر عددی،
چیزی است که این کار می‌خواهد از خود بر جای بگذارد.

\subsection*{آنچه این پژوهش افزود}

سه چیز.

نخست، \textbf{وزن‌دهیِ رژیم‌محور در ترکیب مدل‌ها}. این روش ترکیب ایستا را بهبود
می‌دهد، اما نه به‌طور یکنواخت: تمام بهبود بر روزهای پرتنش متمرکز است و بر
روزهای آرام عملاً هیچ است. همین محدودبودنِ اثر، ادعای قوی‌تری است تا یک بهبودِ
مبهمِ سراسری، چون دقیقاً همان چیزی است که روش برای آن طراحی شده بود. تحلیل
تفسیرپذیری نیز از مسیری مستقل به همان جا رسید و نشان داد رفتار مدل میان دو
رژیم واقعاً تغییر می‌کند.

دوم، \textbf{آزمون خارج از توزیع و تشخیصی که زیر آن نشست}. اجرای مدل‌های
تثبیت‌شده بر داده‌ی امروز، شکستی آشکار را نشان داد؛ اما کار آنجا متوقف نشد.
سنجش پشتیبانیِ ویژگی معلوم کرد که بخش سنجش‌پذیرِ غالبِ آن شکست، ناسازگاریِ
تعریف در مسیر سرویس‌دهیِ خودِ ما بوده است و نه صرفاً دگرگونیِ بازار. تفکیک این
دو، هم نتیجه را دقیق‌تر می‌کند و هم بار مسئولیت را در جای درست می‌گذارد.

سوم، \textbf{یک زنجیره‌ی کامل و بازتولیدپذیر}، از عکس فوریِ داده با اثر انگشت
رمزنگاشتی تا ابزاری که همان مدل تثبیت‌شده را اجرا می‌کند و کیفیت خود را به
کاربر اعلام می‌نماید. مقایسه با نتایج منتشرشده نیز بر همان زنجیره و با کدِ
سنجه‌ی واحد انجام شد، نه با نقلِ اعداد چاپ‌شده.

\subsection*{آنچه این پژوهش نمی‌گوید}

این پژوهش بهترین مدل منتشرشده‌ی مرجع را شکست نمی‌دهد. با دو گونه‌ی آن هم‌تراز
است و به بهترین گونه‌اش می‌بازد. این نتیجه از آغاز پنهان نشد و مسیر دفاع پروژه
پس از ایست‌بازرسیِ میانی آگاهانه از رقابت بر سر جدول امتیازها به سمت مشارکت
روش‌شناختی تغییر کرد.

همچنین نشان داده نشد که یادگیری عمیق در این پیکربندی بر یک مدل خطیِ تنظیم‌شده
مزیت دارد. دو خانواده از یکدیگر تفکیک‌پذیر نبودند، و آنجا که تفاوتی به سود مدل
عمیق دیده شد، تصحیح برای چندگانگی آن را باقی نگذاشت. این خود یافته‌ای است
درباره‌ی آنکه در این مسئله، پیچیدگیِ بیشتر چقدر می‌خرد.

و سرانجام، هیچ‌یک از اعداد دقتِ این پایان‌نامه ادعایی درباره‌ی بازار امروز
نیست. آن‌ها بر دوره‌ای معتبرند که مدل‌ها بر آن آموزش دیده و آزموده شده‌اند، و
\cref{sec:5-2} نشان داد بیرون از آن دوره چه رخ می‌دهد.

% این صفحه یک سطر بلندتر می‌شود تا سطر آخرِ پایان‌نامه تنها روی یک صفحه‌ی
% جداگانه نیفتد؛ \looseness نیز از تِک می‌خواهد بند را یک سطر کوتاه‌تر بچیند.
% متن تغییر نکرده است.
\enlargethispage{2\baselineskip}

\subsection*{جمله‌ی پایانی}

{\looseness=-1
اگر قرار باشد از این کار یک نکته بماند، این است که در پیش‌بینی قیمت برق،
سختیِ اصلی نه در ساختن مدلی دقیق‌تر، که در دانستنِ حدودِ اعتبارِ همان مدل است.
مدل‌های این پژوهش بر دوره‌ی آزمونِ خود خوب عمل کردند و در برخورد با شرایطی که
برای آن ساخته نشده بودند، از یک قاعده‌ی سرانگشتی هم عقب ماندند. آنچه این فاصله
را قابل فهم کرد، نه مدل بهتر بود و نه داده‌ی بیشتر؛ سنجه‌ای بود که پرسید ورودی
امروز اصلاً در دامنه‌ای که مدل دیده است قرار دارد یا نه. پیش‌بینیِ قابل‌اتکا،
پیش از آنکه مسئله‌ی دقت باشد، مسئله‌ی دانستنِ زمانِ بی‌اعتباری است.\par}
